\section{General measures in the minimal envelope class}
\label{sec:general-measure-minimal}

The bounded subsolution theorem and dominated continuity already give the equation for every finite non-pluripolar measure in the minimal envelope class. The missing link is a bounded reference adapted to the measure and to the same background. Its density need only be integrable. We construct that reference from a standard local decomposition theorem, and then give the limiting argument. This reconstructs the minimal-reference case of \cite[Theorem 5.8 and Corollary 5.9]{ALS26}.

\begin{lemma}\label{lem:bounded-measure-carrier}
Let $\theta\geq0$ be smooth and big, and let $\mu$ be a finite non-pluripolar measure on $X$. There are a bounded $\psi\in\PSH(X,\theta)$ and a nonnegative
$f\in L^1(\mathrm{MA}_\theta(\psi))$ such that
\[
 \mu=f\,\mathrm{MA}_\theta(\psi).
\]
No upper or lower volume property is assumed.
\end{lemma}
\begin{proof}
We first specify the local input. Cegrell's decomposition theorem \cite[Theorem 6.3]{Ceg98} says that a positive compactly supported measure $\tau$ on a bounded hyperconvex domain in $\mathbb C^n$, $n\geq2$, is
\[
 \tau=g(\ddc v)^n+\tau_s,
 \qquad v\in\mathcal E_0,\qquad
 0\leq g\in L^1((\ddc v)^n),
\]
where $\tau_s$ is carried by a pluripolar set. Here the local class $\mathcal E_0$ consists of bounded negative plurisubharmonic functions with zero boundary values and finite total Monge--Amp\`ere mass. This local notation is different from the global envelope class $\mathcal E_0$ in this chapter. If $\tau$ is non-pluripolar, $\tau_s=0$, so $\tau\ll(\ddc v)^n$. A change in the normalization of $\ddc$ only changes $g$ by a positive constant.

For completeness the required local assertion on a disk also follows from the Green potential. Let $a\leq0$ be the Green potential of a finite compactly supported non-polar measure $\tau$, normalized by $\ddc a=\tau$. Then $v=e^a$ is bounded and subharmonic, and
\[
 \ddc v\geq e^a\tau.
\]
To justify the inequality, set $a_M=\max(a,-M)$. The bounded subharmonic chain rule gives
$\ddc e^{a_M}\geq e^{a_M}\ddc a_M$, which on $\{a>-M\}$ is at least $e^a\tau$. Let $M\to\infty$, using decreasing convergence of the bounded subharmonic functions $e^{a_M}$ and monotone convergence on the right. Since $\{a=-\infty\}$ is polar, $e^a>0$ almost everywhere for $\tau$. Thus $\tau\ll\ddc v$. The bounded chain inequality itself follows by local convolution and passage to the distributional limit. We use only this absolute continuity assertion in every dimension.

Choose a strict potential $\rho$ with analytic singularities, smooth off its analytic pole set $Z$, and
\[
 \theta+\ddc\rho\geq\kappa\omega,\qquad\kappa>0.
\]
Such a potential is available by strict regularization. Cover $X\setminus Z$ by countably many compact inner coordinate balls $K_i$, with larger coordinate balls $B_i\Subset X\setminus Z$. Apply the local assertion to $\mu_i=\mathbf1_{K_i}\mu$ on $B_i$. After a positive affine change of its bounded local potential, we obtain $v_i\in\PSH(B_i)$ with
\[
 1\leq v_i\leq2,\qquad
 \mu_i\ll(\ddc v_i)^n.
\]
Zero pieces may be omitted. Pick a smooth $\chi_i$ on $B_i$ which is zero near $K_i$ and $-3$ near $\partial B_i$. For a sufficiently small $\epsilon_i>0$,
\[
 \theta+\ddc\bigl(\rho+\epsilon_i(v_i+\chi_i)\bigr)
 \geq\tfrac\kappa2\omega+\epsilon_i\ddc v_i.
\]
Inside $B_i$ take the maximum of this potential with $\rho$, and extend it by $\rho$ outside $B_i$. The two pieces glue because $v_i+\chi_i\leq-1$ on a boundary collar. Finally take the maximum with $-C_i$, where $C_i$ is large enough that this last maximum does not change the function near $K_i$. Constants are $\theta$-plurisubharmonic since $\theta\geq0$. We have obtained a global bounded $q_i\in\PSH(X,\theta)$ which agrees with $\rho+\epsilon_i v_i$ near $K_i$. Positivity of bounded mixed products therefore gives
\[
 \mathrm{MA}_\theta(q_i)\geq\epsilon_i^n(\ddc v_i)^n
 \quad\hbox{near }K_i,
 \qquad \mu_i\ll\mathrm{MA}_\theta(q_i).
\]

Translate each $q_i$ so that $q_i\leq0$, and choose positive numbers $t_i$ with
\[
 \sum_i t_i<1,\qquad
 \sum_i t_i\|q_i\|_\infty<\infty.
\]
The uniformly convergent sum $\psi=\sum_i t_iq_i$ is bounded and $\theta$-plurisubharmonic. For a fixed index $i$ it is a convex combination
$\psi=t_iq_i+(1-t_i)w_i$ with a bounded $w_i\in\PSH(X,\theta)$: use the remaining terms and the unused weight on the constant zero potential. Hence
\[
 \mathrm{MA}_\theta(\psi)\geq t_i^n\mathrm{MA}_\theta(q_i).
\]
Every $\mu_i$ is absolutely continuous with respect to this single finite measure. The sets $K_i$ cover $X\setminus Z$, and $\mu(Z)=0$. Thus $\mu\ll\mathrm{MA}_\theta(\psi)$. The Radon--Nikodym theorem supplies the asserted integrable density. If $\mu=0$, simply take $\psi=0$ and $f=0$.
\end{proof}

\Needspace{15\baselineskip}
\begin{theorem}\label{thm:general-measure-minimal}
Let $\theta\geq0$ be smooth and big, and assume $v_-(\theta)>0$. Let $\mu$ be a finite non-pluripolar measure with $\mu(X)>0$.
For every $\lambda>0$ there is a unique $u\in\mathcal E_0$ satisfying
\[
 \mathrm{MA}_\theta^{\mathrm{ext}}(u)=e^{\lambda u}\mu.
\]
There also exist $v\in\mathcal E_0$, normalized by $\sup_Xv=0$, and a unique positive scalar $c$ such that
\[
 \mathrm{MA}_\theta^{\mathrm{ext}}(v)=c\mu.
\]
The second assertion does not claim uniqueness of the normalized potential. No bounded mass property is assumed.
\end{theorem}
\begin{proof}
The reference $0$ is model because $\theta\geq0$, and $P_{\epsilon\theta}(0)=0$ for every $\epsilon>0$. Every bounded $\theta$-plurisubharmonic function belongs to $\mathcal E_0$, since a sufficiently negative constant is a competitor for each of its stretched envelopes. The lower estimate of \cref{thm:relative-uniform-mass-lower}, with $\epsilon\downarrow0$, gives
\begin{equation}\label{eq:minimal-lower-mass}
 \mathrm{MA}^{\mathrm{ext}}_\theta(w)(X)\geq v_-(\theta)
 \qquad(w\in\mathcal E_0).
\end{equation}

\emph{Positive exponent.} Write $\mu=f\nu$ by the preceding lemma, with $\nu=\mathrm{MA}_\theta(\psi)$ and $\psi$ bounded. Let $f_j=\min(f,j)$ for positive integers $j$. These densities are bounded and nonzero as measures. The bounded subsolution theorem gives bounded $u_j$ with
\[
 \mathrm{MA}_\theta(u_j)=e^{\lambda u_j}f_j\nu.
\]
Exponential order shows that $u_j$ decreases, because $f_j$ increases. In particular $u_j\leq u_1$. On the other hand
\[
 v_-(\theta)\leq\mathrm{MA}_\theta(u_j)(X)
       \leq e^{\lambda\sup_Xu_j}\mu(X).
\]
Thus their suprema have a common lower bound. The decreasing limit $u$ is a genuine $\theta$-plurisubharmonic function, and convergence holds in $L^1(X)$. Moreover
\[
 \mathrm{MA}_\theta(u_j)\leq e^{\lambda\sup_Xu_1}\mu.
\]
Apply \cref{thm:relative-fixed-measure-continuity} with reference zero. It gives $u\in\mathcal E_0$ and weak convergence of the Monge--Amp\`ere measures to $\mathrm{MA}^{\mathrm{ext}}_\theta(u)$. The right sides converge in total variation to $e^{\lambda u}f\nu$: the functions $e^{\lambda u_j}f_j$ converge almost everywhere for $\nu$ and are bounded by $e^{\lambda\sup_Xu_1}f\in L^1(\nu)$. This proves the equation. Its measure is finite, so \cref{cor:relative-finite-exponential-order} gives uniqueness.

\Needspace{10\baselineskip}
\emph{Zero exponent.} For each positive integer $j$, let $u_j\in\mathcal E_0$ be the just-constructed solution with exponent $1/j$. We first bound the normalization. Let $h$ be the bounded solution of
\[
 \mathrm{MA}_\theta(h)=e^h\min(f,1)\nu,
\]
supplied by the bounded subsolution theorem. Set $H=\|h\|_\infty$ and $w_j=h+(j-1)H$. Then $w_j/j\geq h$, and
\[
 e^{-w_j/j}\mathrm{MA}_\theta(w_j)
 \leq e^{-h}\mathrm{MA}_\theta(h)
 \leq\mu
 =e^{-u_j/j}\mathrm{MA}^{\mathrm{ext}}_\theta(u_j).
\]
Exponential order, now with exponent $1/j$, gives $u_j\leq w_j$. Put
\[
 M_j=\sup_Xu_j,\qquad v_j=u_j-M_j,
 \qquad c_j=e^{M_j/j}.
\]
Then $\sup_Xv_j=0$ and $0<c_j\leq e^H$. Equation \eqref{eq:minimal-lower-mass} also gives
\[
 v_-(\theta)\leq c_j\int_X e^{v_j/j}\,\mathrm{d}\mu
                 \leq c_j\mu(X),
\]
so $c_j\geq v_-(\theta)/\mu(X)>0$. After extraction, compactness of normalized potentials gives $v_j\to v$ in $L^1(X)$, with $v$ genuine and normalized, and $c_j\to c>0$.

By \cref{lem:normalized-exponential-npp}, $a_j=e^{v_j/j}\to1$ in $L^1(\mu)$. Consequently
\[
 \mathrm{MA}^{\mathrm{ext}}_\theta(v_j)=c_ja_j\mu
       \longrightarrow c\mu
 \quad\hbox{in total variation},\qquad
 \mathrm{MA}^{\mathrm{ext}}_\theta(v_j)\leq e^H\mu.
\]
Apply \cref{thm:relative-fixed-measure-continuity} again. It gives $v\in\mathcal E_0$ and identifies the limit as $\mathrm{MA}^{\mathrm{ext}}_\theta(v)$. The scalar assertion in \cref{cor:relative-finite-exponential-order}, applied in both directions to two solutions, proves uniqueness of $c$.
\end{proof}

Both equations in the minimal class are obtained using domination by one fixed finite measure. This result also provides comparison potentials for a prescribed reference: finite-measure exponential order can be applied with reference zero and any second potential more singular than zero. The proof in \cref{sec:general-measure-relative} uses this observation after its first approximation with moving bounded reference measures.
